\documentclass{cssheet}

%--------------------------------------------------------------------------------------------------------------
% Basic meta data
%--------------------------------------------------------------------------------------------------------------

\title{Mathematische Strukturen in Okkultismus, Esoterik und Religion}
\author{Prof. Dr. Christian Spannagel}
\date{\today}
\hypersetup{%
    pdfauthor={\theauthor},%
    pdftitle={\thetitle},%
    pdfsubject={Themen Masterarbeit},%
    pdfkeywords={master, phhd}
}

%--------------------------------------------------------------------------------------------------------------
% document
%--------------------------------------------------------------------------------------------------------------

\begin{document}

\vspace*{5mm}
\begin{center}
{\Large Themen für eine Masterarbeit}
\end{center}

\printtitle
\vspace*{1cm}

Okkultismus und Magie sind seit jeher eng mit der Mathematik verknüpft. Magische Quadrate, geometrische Figuren wie das Pentagramm oder Zahlen in der Numerologie wurden mit besonderen Bedeutungen aufgeladen. Auch in religiösen Kontexten schreibt man bestimmten Zahlen und geometrischen Formen besondere Eigenschaften zu.

Jenseits religiöser, okkulter und esoterischer Deutungen bieten viele dieser Symbole spannende Anknüpfungspunkte für die Auseinandersetzung mit mathematischen Strukturen. Ihre Behandlung im Mathematikunterricht erfordert zugleich einen fachlich fundierten und pädagogisch angemessenen Umgang mit den damit verbundenen religiösen, weltanschaulichen und esoterischen Vorstellungen. Religiöse und weltanschauliche Bezüge sollen sachlich, respektvoll und differenziert eingeordnet werden. Esoterische beziehungsweise nichtwissenschaftliche Aussagen sollen anhand wissenschaftlicher Kriterien kritisch reflektiert werden.

Im Rahmen einer Masterarbeit soll ein ausgewähltes Themengebiet sowohl mathematisch als auch mathematikdidaktisch untersucht werden. Hierzu zählen z.\,B. die folgenden Themen oder Teilbereiche:

\begin{itemize}
\item Sternpolygone (Pentagramm, Hexagramm, Heptagramm)
\item Gematrie und Zahlen in der Bibel
\item Symbole der Unendlichkeit
\item Mathematik im Volksglauben
\item Magische Quadrate
\item Heilige Geometrie
\item Numerologie und Mathematik
\item Kryptographie in okkulten Kontexten
\item Mathematik in Horoskopen
\item Tarot
\end{itemize}

Eigene Schwerpunktthemen können gerne vorgeschlagen werden.

Die theoretische und mathematikdidaktische Analyse soll in die Konzeption einer Unterrichtseinheit einfließen. Diese wird in der Schule durchgeführt und im Rahmen der Masterarbeit empirisch untersucht.

Mögliche Tätigkeiten sind:

\begin{itemize}
\item \textbf{Literaturrecherche:} Untersuchung der kulturhistorischen Bedeutung und der mathematischen Besonderheiten des gewählten Symbols, Gegenstands oder Zahlenbereichs.

\item \textbf{Analyse:} Mathematische und mathematikdidaktische Analyse der zugrunde liegenden Strukturen sowie differenzierte Einordnung religiöser und weltanschaulicher Bezüge und kritische Auseinandersetzung mit esoterischen beziehungsweise nichtwissenschaftlichen Aussagen.

\item \textbf{Entwicklung:} Konzeption einer Unterrichtseinheit zur Behandlung des Themenbereichs in der Schule. Dabei sollen die mathematischen Inhalte fachlich fundiert und die religiösen, weltanschaulichen und/oder esoterischen Bezüge pädagogisch angemessen berücksichtigt werden.

\item \textbf{Durchführung und empirische Untersuchung:} Erprobung der Unterrichtseinheit in der Schule sowie empirische Untersuchung ausgewählter Aspekte, beispielsweise der Lernprozesse, der Vorstellungen und Argumentationen der Schülerinnen und Schüler oder der Wahrnehmung und Bewertung der Unterrichtseinheit.

\item \textbf{Evaluation:} Methodisch nachvollziehbare Auswertung und Diskussion der erhobenen Daten sowie Ableitung von Konsequenzen für die Weiterentwicklung der Unterrichtseinheit.
\end{itemize}

\vspace*{10mm}

\printlicense

\end{document}
